\documentclass[../../../include/open-logic-section]{subfiles}

\begin{document}

\olfileid[pt]{sth}{card-arithmetic}{opps}
\olsection{Definindo as Operações Básicas}

Como não precisamos de controlar a ordem, a aritmética cardinal é
bastante mais fácil de definir do que a aritmética ordinal. Definiremos
adição, multiplicação e exponenciação simultaneamente. 

\begin{defn}
Quando $\cardfont{a}$ e $\cardfont{b}$ são cardinais:
\begin{align*}
	\cardfont{a} \cardplus \cardfont{b} &\defis 
	\card{\cardfont{a} \disjointsum \cardfont{b}}\\
	\cardfont{a} \cardtimes \cardfont{b} &\defis 
	\card{\cardfont{a} \times \cardfont{b}}\\
	\cardexpo{\cardfont{a}}{\cardfont{b}} &\defis 
	\card{\funfromto{\cardfont{b}}{\cardfont{a}}}
\end{align*}
onde $\funfromto{X}{Y} = \Setabs{f}{f \text{ é uma função } X \to Y}$.
(É fácil mostrar que $\funfromto{X}{Y}$ existe para quaisquer conjuntos $X$
e $Y$; deixamos isto como exercício.) 
\end{defn}

\begin{prob}
Prove em $\Zminus$ que $\funfromto{X}{Y}$ existe para quaisquer conjuntos $X$
e~$Y$. Trabalhando em $\ZF$, calcule $\setrank{\funfromto{X}{Y}}$ a partir de
$\setrank{X}$ e $\setrank{Y}$, à maneira de
\olref[sth][ord-arithmetic][using-addition]{rankcomputation}. 
\end{prob}

Pode ajudar explicar esta definição. Quanto à adição: usa-se a noção de soma disjunta, $\disjointsum$, tal como definida em
\olref[ord-arithmetic][add]{defdissum}; e é fácil
ver que esta definição dá o resultado certo nos casos finitos.
Quanto à multiplicação: \olref[sfr][set][pai]{cardnmprod} nos diz
que, se $A$ tem $n$ membros e $B$ tem $m$ membros, então $A \times B$
tem $n \cdot m$ membros, pelo que a nossa definição generaliza simplesmente a ideia
à multiplicação transfinita. A exponenciação é semelhante: estamos
simplesmente a generalizar o pensamento do finito para o transfinito.
De fato, sob certos aspectos, a aritmética cardinal transfinita parece muito
mais a aritmética ``ordinária'' do que a aritmética ordinal
transfinita:

\begin{prop}\ollabel{cardplustimescommute}
$\cardplus$ e $\cardtimes$ são comutativas e associativas. 
\end{prop}

\begin{proof}
Para a comutatividade, por
\olref[cardinals][cardsasords]{lem:CardinalsBehaveRight} basta
observar que $\cardeq{(\cardfont{a} \disjointsum
\cardfont{b})}{(\cardfont{b} \disjointsum \cardfont{a})}$ e
$\cardeq{(\cardfont{a} \times \cardfont{b})}{(\cardfont{b} \times
\cardfont{a})}$. Deixamos a associatividade como exercício.
\end{proof}

\begin{prob}
Prove que $\cardplus$ e $\cardtimes$ são associativas.
\end{prob}

\begin{prop}
$A$ é infinito se e só se $\card{A} \cardplus 1 = 1 \cardplus \card{A} = \card{A}$.
\end{prop}

\begin{proof}
Como em
\olref[cardinals][classing]{generalinfinitycharacter}, a partir de
\olref[ord-arithmetic][using-addition]{ordinfinitycharacter} e
\olref[cardinals][cardsasords]{lem:CardinalsBehaveRight}. 
\end{proof}

Isto explica porque precisamos de símbolos distintos para adição/multiplicação ordinal versus
cardinal: são operações genuinamente \emph{diferentes}.
Este par de resultados mostra que a exponenciação ordinal versus
cardinal também são operações distintas. (Recordar que
\olref[z][infinity-again]{defnomega} implica que $2 = \{0,
1\}$):

\begin{lem}\ollabel{lem:SizePowerset2Exp}
$\card{\Pow{A}} = \cardexpo{2}{\card{A}}$, para qualquer $A$.
\end{lem}

\begin{proof}
Para cada subconjunto $B \subseteq A$, seja $\chi_B \in \funfromto{A}{2}$ dado por:
\begin{align*}
	\chi_{B}(x) &\defis
	\begin{cases}
		1 & \text{se }x\in B\\
		0 & \text{caso contrário.}
	\end{cases}
\end{align*}
Seja agora $f(B) = \chi_B$; isto define !!a{bijection} $f \colon \Pow{A}
\to \funfromto{A}{2}$. Logo $\cardeq{\Pow{A}}{\funfromto{A}{2}}$. Daí
$\cardeq{\Pow{A}}{\funfromto{\card{A}}{2}}$, pelo que
$\card{\Pow{A}} = \card{\funfromto{\card{A}}{2}} =
2^{\card{A}}$.
\end{proof}

Esta prova sucinta subsume essencialmente a discussão em
\olref[sfr][siz][red-alt]{sec}. Ali, mostramos como ``reduzir'' a
não enumerabilidade de $\Pow{\omega}$ à não enumerabilidade do conjunto das
sequências binárias infinitas, $\Bin^\omega$. Na prática, $\Bin^{\omega}$ é
apenas $\funfromto{\omega}{2}$; e a prova precedente mostrou que o
raciocínio que seguimos na \olref[sfr][siz][red-alt]{sec} continua a
valer usando qualquer conjunto~$A$ no lugar de~$\omega$. O resultado também fornece
um fato rápido sobre exponenciação cardinal:

\begin{cor}\ollabel{cantorcor}
$\cardfont{a} < \cardexpo{2}{\cardfont{a}}$ para qualquer cardinal~$\cardfont{a}$.
\end{cor}

\begin{proof}
Pelo Teorema de Cantor (\olref[sfr][siz][car]{thm:cantor}) e
\olref{lem:SizePowerset2Exp}.
\end{proof}
\noindent
Logo $\omega < \cardexpo{2}{\omega}$. Mas nota: isto é um resultado sobre
exponenciação \emph{cardinal}. Deve contrastar-se com a
exponenciação \emph{ordinal}, pois neste último caso $\omega =
\ordexpo{2}{\omega}$ (ver \olref[ord-arithmetic][expo]{sec}).

Enquanto estamos no tema da exponenciação cardinal, podemos também ser
um pouco mais precisos sobre a ``maneira'' pela qual $\Real$ é
!!{nonenumerable}.

\begin{thm}\ollabel{continuumis2aleph0}
$\card{\Real} = \cardexpo{2}{\omega}$
\end{thm}

\begin{proof}[Esboço de prova]
Há muitas maneiras de provar isto. A mais direta é
argumentar que $\cardle{\Pow{\omega}}{\Real}$ e
$\cardle{\Real}{\Pow{\omega}}$, e depois usar Schr\"oder-Bernstein para
inferir que $\cardeq{\Real}{\Pow{\omega}}$, e
\olref[card-arithmetic][opps]{lem:SizePowerset2Exp} para inferir
que $\card{\Real} = \cardexpo{2}{\omega}$. Deixamos como
exercício (esclarecedor) definir injecções $f \colon
\Pow{\omega} \to \Real$ e $g \colon \Real \to \Pow{\omega}$.
\end{proof}

\begin{prob}
Complete a prova de
\olref[sth][card-arithmetic][opps]{continuumis2aleph0},
mostrando que $\cardle{\Pow{\omega}}{\Real}$ e
$\cardle{\Real}{\Pow{\omega}}$.
\end{prob}

\end{document}
